% Capítulo 11. Conclusiones (rúbrica 25 pts).

\section{Hallazgos}
\label{sec:hallazgos}

El problema de programar la flota eléctrica de RED se pudo resolver para toda la red y los tres tipos de día. El plan obtenido cubre las mismas \CorridasLaboral{} corridas del día laboral que un caso base razonable con \AhorroBusesLaboral{} buses menos (\DeltaPctBusesLaboral\,\%) y un costo \BajaCostoLaboral\,\% menor, y en la semana ahorra cerca de USD~\AhorroSemanaMillones{} millones (\DeltaPctSemana\,\%). Ese era el objetivo que se planteó al comienzo, operar el mismo servicio con menos flota.

El ahorro viene casi todo de los buses, que explican cerca de tres de cada cuatro dólares, y la mayor parte se consigue antes de optimizar. Compartir buses entre las rutas de un sector, cargar solo lo necesario y elegir el bus más cercano ya entrega el \PctAhorroInicial\,\% de la reducción de flota. La búsqueda por grupos agrega el resto, y entre las variantes probadas la mejor fue tomar cuatro buses al azar, que es el punto donde el tamaño del grupo y la cantidad de intentos se equilibran.

Durante el día la infraestructura de carga sobra. Como máximo se usa cerca de un quinto de las conexiones de un electroterminal, las plazas solo se llenan al cierre y ampliar o reducir la capacidad casi no cambia la flota. De noche ocurre lo contrario. Para repetir el día laboral habría que reponer \ReponerGwhLaboral~GWh, y las conexiones actuales solo alcanzan para el \PctCubiertoLaboral\,\% de esa energía. Esa es la restricción que viene para RED, y afecta también al caso base.

La solución es aproximada. En los casos donde el óptimo se pudo demostrar queda en promedio a \BrechaCuartetos\,\% de él, en un caso mediano llega a \BrechaCuartetosRutasOchoAbs\,\%, y cuando se obliga a recargar queda cerca de un \DifMediaRecargaCientoCincuentaAbs\,\% sobre el mejor plan que encontró el modelo exacto, al que en el caso más grande supera. En la red completa solo se puede acotar, y queda \BrechaCotaCosto\,\% sobre una cota débil.

\section{Fortalezas y limitaciones}
\label{sec:fortalezas}

La principal fortaleza del trabajo es que el método escala a la red completa sin perder control sobre la factibilidad. Cada plan pasa por un auditor independiente, la búsqueda solo acepta cambios factibles y más baratos, y las cifras del informe se generan desde los resultados, de modo que se pueden rastrear y repetir. La segunda es que la calidad del método se midió con el modelo exacto en casos chicos y medianos. La tercera es que la conclusión principal resiste los cambios de precios, de capacidad de carga y del bus, y aun cobrando la recarga nocturna el ahorro supera el 12\,\%.

Las limitaciones más importantes son dos. El plan es más sensible a los atrasos que el caso base, porque tiene menos buses y carga lo justo, y con tiempos de viaje 10\,\% variables el \AtrasadasPlanDiez\,\% de sus corridas sale atrasado más de cinco minutos. Y el modelo trata cada día por separado, sin exigir que la energía al cierre permita repetirlo. A eso se suman los sectores fijos, que dejan fuera combinaciones entre sectores vecinos, los supuestos sobre la velocidad en vacío y la carga lineal, la ausencia de conductores y mantenciones en el modelo, y una cota débil para la red completa. Mirando hacia atrás, habría convenido incorporar antes la noche y la holgura al modelo en lugar de medirlas al final.

\section{Pasos siguientes}
\label{sec:pasos}

El paso más urgente es agregar holgura de energía y de tiempo al MILP de recargas y medir cuánto ahorro se pierde con ella, porque sin eso el plan no se puede operar. Le siguen modelar la noche y ampliar las operaciones de la búsqueda. Con esas mejoras, la recomendación para RED es un piloto en el sector de Vespucio Norte, primero en paralelo y luego operado frente a un sector de control, con criterios de éxito definidos de antemano en buses, puntualidad y energía, y una evaluación de la ampliación de la carga nocturna antes de extender el plan a toda la red.
